\chapter{The GLoRa Contribution}
\label{chap:glora-contribution}

A benchmark score supports a claim about dependency length only if success requires the model to learn that dependency.
The included paper introduces GLoRa, a benchmark generator for testing whether a graph learning system can learn dependencies of a specified length \cite{zhou2025glora}.
For a chosen dependency-length parameter $d$, GLoRa generates balanced node-classification examples whose labels depend on a source-to-target path condition.
It then controls the construction so that fitting sufficiently many examples requires the intended dependency rather than a shortcut function.

The construction addresses six threats to this interpretation.
The path-aware target and Property (P1) address path dependence, while alternative chains and additional holes address hole-count shortcuts and graph statistics.
Property (P2) ensures that the target remains expressible by the evaluated GNN-based systems, and Property (P3) addresses fairness between the two classes.
The directed construction and the representation and gradient measurements reduce diagnostic ambiguity when performance falls.
The experiments then use accuracy over increasing $d$ to estimate where each system stops learning a fitting function under these controls.

\section{Motivation}
\label{sec:glora-motivation}

Architectural proposals target over-smoothing, over-squashing, and vanishing gradients, but comparing them requires a benchmark that identifies what successful prediction establishes \cite{li2018deeper,alon2021bottleneck,glorot2010difficulty}.

The first evaluation threat is a failure to establish path dependence.
For dependency length $d$, the relevant question is whether the target prediction depends on the distance and path information of a source-to-target path of length $d$.
This requirement is stronger than high accuracy on labels generated by a long-range rule.
A dataset can admit more than one fitting function.
One function may use the intended source-to-target path of length $d$.
Another shortcut function may exploit a local count, a global graph statistic, a source-node marker, or an artefact of how positive and negative examples were generated.
If both functions fit the examples, high test accuracy cannot distinguish a model that learned the path-aware dependency from a model that learned the shortcut.

Real-world benchmarks compare practical performance, but they do not formally control the dependency that a learned function must use \cite{dwivedi2022long,zhou2025glora}.

Synthetic benchmarks offer more control, but that control must extend to alternative fitting functions.
A synthetic benchmark can choose the graph structure, the node features, and the target function.
This makes it possible to set the dependency length directly.
Two shortcut threats remain.
If every positive example has no hole and every negative example has one visible hole, a model can solve the task by detecting the hole rather than by testing whether information is propagated from the source to the target through a valid path.
If the only source node has a unique local degree pattern, a model can detect the source rather than learning the source-to-target path condition.
If graph size differs between classes, a graph-level statistic may solve the benchmark.
The generator must therefore address both hole-count shortcuts and predictive graph statistics if it is to make the intended path-aware dependency necessary rather than merely present.

GLoRa turns these requirements into three properties developed later in the chapter.
Fitting sufficiently many generated examples should require a dependency of length $d$.
At least one fitting function should be expressible by the GNN-based systems under evaluation, and positive and negative examples should come from the same stochastic generator \cite{zhou2025glora}.
The last condition addresses fairness by avoiding an artificial distribution shift between classes.
Together, the conditions define a benchmark generator rather than a fixed dataset.
For each generator parameter $d$, it creates a benchmark whose role is to test learning at that length.
Running a model over increasing values of $d$ then gives an empirical estimate of the largest dependency length that the model can learn under the benchmark conditions.

\section{Path-Aware Dependency and Notation}
\label{sec:glora-path-aware}

GLoRa addresses the path-dependence threat through a path-aware target.
The target label depends on a source-to-target path of length $d$, not merely on the presence of a distant source node.
The notation below fixes the graph, path, source, and target so that later claims about dependency length state exactly what evidence the benchmark certifies.

The paper works with directed graphs, with an undirected variant obtained by symmetrising the edges.
Throughout this chapter, a graph is written as
\[
  G=(V,E,\lambda),
\]
where $V$ is the finite set of nodes, $E \subseteq V \times V$ is the set of directed edges, and $\lambda$ maps each node to its feature vector.
For a dependency length $d$, write the relevant source-to-target path as
\[
  p=(v_0,\ldots,v_d).
\]
The source node is $v_0$ and the target node is $v_d$.
The intermediate nodes are $v_1,\ldots,v_{d-1}$.
The edges follow the path direction, so $(v_{i-1},v_i) \in E$ for every $i \in \{1,\ldots,d\}$.
The target prediction is the label assigned to $v_d$.

The paper's path-aware formalization is deliberately narrow.
It does not try to capture every possible meaning of long-range dependence in graphs.
Instead, it captures a dependency where the target label depends on information that must be preserved along a particular source-to-target path.
In this setting, it is not enough that a source node at large graph distance exists somewhere in the graph.
The prediction must depend on the relation between the source $v_0$, the target $v_d$, and the intermediate nodes that connect them.

The formal definition in the paper can be read as three tests on a node-classification function.
First, each intermediate node on the witness path must matter.
If the feature of any $v_i$ with $0<i<d$ is changed in a designated way, the output at $v_d$ changes.
Second, the function must keep the same value when a duplicate of part of the path is changed but the original path remains intact.
This prevents the definition from accepting functions that respond merely to the presence of a local feature value somewhere in the graph.
Third, if both the original path and the duplicate are changed, the output changes again.
Together, these requirements say that the source-to-target path is not decoration around a simpler feature test.
The prediction depends on the way information is arranged along the path from $v_0$ to $v_d$.

This formalization states what the generator must make necessary.
The benchmark should generate enough examples that any function fitting them must rely on a path-aware dependency of length $d$.
This property distinguishes GLoRa from a benchmark whose labels are produced by a long-range rule but can also be fitted by a short-range shortcut function.
The paper also introduces an approximate version of the definition, parameterised by a precision $\delta$.
This is needed because GLoRa uses real-valued node features sampled from intervals rather than a fixed finite alphabet.
The approximate definition preserves the same intuition: perturbing the relevant path features within a controlled tolerance should witness the dependency.

\section{Benchmark Design}
\label{sec:glora-design}

The GLoRa generator turns the path-aware definition into balanced inductive node-classification examples.
Positive and negative graphs differ by whether a valid source-to-target path exists.
Each remaining design element answers a specific threat to interpreting accuracy: alternative chains address hole-presence shortcuts, additional holes address hole-count and related graph statistics, and directed edges control the source-to-target path count for later diagnosis.
Together, these controls make dependency length the experimental variable rather than an incidental property of the generated graphs.

GLoRa generates inductive binary node-classification examples.
Each example consists of a graph $G=(V,E,\lambda)$ and a target node.
The label says whether the target node is positive or negative.
The graphs use two-dimensional node features.
The first coordinate marks a source node, and the second coordinate marks whether a node can participate in a valid path.
For exposition, it is useful to describe three node types:
\begin{description}
\item[Source.]
The source node $v_0$ has a feature of the form $(1,y)$.
The first coordinate marks the node as a source, and the second coordinate is sampled from a range that avoids the special values used by the path condition.
\item[Normal path node.]
A normal non-source node has a feature of the form $(x,1)$.
The second coordinate marks it as usable on a valid path.
\item[Hole.]
A hole has a feature of the form $(x,0)$.
It breaks the intended path condition.
\end{description}
The remaining coordinate values are written here as $x$ and $y$.
They are sampled from a bounded set that stays away from $0$ and $1$.
In the paper, this set is $[-12,-2]\cup[3,13]$.
This choice keeps the special marker values separated from the random filler values \cite{zhou2025glora}.

The feature schema defines a simple path-aware target function.
The target node $v_d$ is positive if and only if there is a source-to-target path from a source node $v_0$ to $v_d$ such that all non-source nodes on the path have second coordinate $1$.
In the idealised picture, the witness path is $p=(v_0,\ldots,v_d)$.
Every node $v_i$ for $0<i\leq d$ has a normal path feature $(x_i,1)$.
If any intermediate node becomes a hole, the source-to-target path no longer witnesses a positive label.
The chapter uses length in two distinct but related ways.
In the thesis notation, an exact witness path $p=(v_0,\ldots,v_d)$ has length $d$.
In the implemented GLoRa generator, the parameter $d$ instead controls the scale of the sampled source-to-target path, whose exact length may vary near that value so that path length itself is not a class marker.
Experimental settings reported by $d$ refer to this generator parameter unless the text explicitly defines an exact witness path.

The first part of the generator creates the main source-to-target path and establishes the class-defining path condition.
For a positive example, this path has no hole.
For a negative example, one non-target intermediate node on the main source-to-target path is turned into a hole.
If the generator stopped here, it would not address the hole-presence shortcut.
A classifier could simply ask whether the graph contains any hole outside the target node.
That question does not require using distance and path information from the source to the target.

The second part adds alternative chains to remove this hole-presence shortcut.
Each alternative chain connects two points that already lie on previously generated chains.
Each such chain contains a hole.
These chains make holes common in both positive and negative examples.
After this step, the existence of a hole is no longer enough to distinguish the labels.
The classifier must identify whether there is at least one complete source-to-target path whose non-source nodes are all normal.
It must reason about the source-to-target path condition, not only about the multiset of node features.

The third part addresses graph-statistic shortcuts by adding further holes.
In particular, this step equalises the distribution of hole counts between positive and negative examples.
For positive examples, the extra holes are placed away from the main source-to-target path so that it remains valid.
For negative examples, the main source-to-target path already contains a hole.
The added holes make it harder for a classifier to solve the task by counting holes or by using a rule such as ``positive graphs have one fewer hole than negative graphs.''
The generator therefore combines a clear intended dependency with randomised distractors that remove the most obvious shortcut functions.

Directedness addresses a later source of diagnostic ambiguity.
Because edges point toward the target, the number of paths leading into $v_d$ is controlled by the number of additional chains.
This number is sampled between five and ten in the paper's algorithm.
The bound does not grow with $d$.
This bound lets the diagnostic argument separate increasing dependency length from the path proliferation associated with over-squashing.
In an undirected version, even a simple path can create many walks when message passing can use both edge directions.
The directed construction keeps the evaluation focused on longer dependency length rather than on an exponentially growing number of message-passing walks.

\section{Examples and the Role of Shortcuts}
\label{sec:glora-shortcuts}

Small GLoRa examples show why a generated long-range rule does not by itself establish path dependence.
A clean source-to-target path makes the target positive, but the benchmark must stop a model from replacing that condition with a cheaper rule such as counting holes.
Alternative chains and added holes are designed to rule out these explanations for a high score.

Consider a benchmark containing only the main chains described above.
Its positive example has a source $v_0$, a target $v_d$, and a path $p=(v_0,\ldots,v_d)$ between them.
The source has first coordinate $1$, and every non-source node on the path has second coordinate $1$.
The target is therefore positive.
A negative example differs only in that one intermediate node $v_i$ has second coordinate $0$, making it a hole that breaks the source-to-target path.
Although this rule is long-range, these examples do not force the learned function to use it.
A model with a global readout, a virtual node, or a preprocessing step that exposes graph-level statistics could classify the target as positive whenever the graph contains no hole.
This shortcut fits the data without using the path from $v_0$ to $v_d$, because neither the location of the hole nor whether it lies between the source and target affects the prediction.
The function would learn an artefact of the example construction rather than the intended dependency.

The graph-statistic threat extends beyond hole counts.
If the source is always the only node with no incoming edges, a model may detect the source by local degree rather than by a path condition.
If positive examples have a different number of nodes or edges than negative examples, graph size can become predictive.
If additional paths appear only in one class, the number of paths can become a class marker.
If the target node has a different graph neighbourhood in positive and negative examples, a short-range message-passing model can perform well without using information from nodes at large graph distance.
These shortcut functions are not implementation details.
They determine what conclusion the benchmark score can support.

Appendix~D of the paper uses this lens to examine existing synthetic benchmarks \cite{zhou2025glora}.
The criticism is not that those benchmarks are badly designed for every purpose.
The point is that they do not certify the particular claim GLoRa is designed to test.
Synthetic Chains, Graph Transfer, Ring Transfer, and Tree-Neighbours all contain long-range generative rules, but they also admit fitting functions that do not rely on the intended path-aware dependency.
h-Proximity has a related issue: a fitting function can avoid depending on the features along the paths between relevant nodes and the target.
Colour-Connectivity has the opposite problem.
It does require a global dependency, but the paper argues that the relevant fitting function is not expressible by the GNN-based systems under consideration.
Conditional Recall and Tree-Max use graph families with limited size, which again makes it possible to solve the task without certifying long-range dependency learning at arbitrary length.

The shortcut analysis establishes what the generator must make necessary, but necessity alone is not enough.
Alternative chains and holes in both classes require a fitting function to determine whether a valid source-to-target path exists under the feature pattern.
The benchmark must also keep that function expressible by standard GNN-based systems, or a failure would reflect an impossible target rather than an inability to learn the intended dependency.
Its difficulty is therefore targeted rather than arising merely because the examples are random, large, or noisy.

\section{Why Expressibility Matters}
\label{sec:glora-expressibility}

A second evaluation threat is that the target function may be inexpressible for the evaluated model class.
GLoRa therefore treats expressibility as part of benchmark design.
Under this control, failure concerns learning the intended path-aware dependency rather than inexpressibility.

The paper states three intended properties that divide these evaluation responsibilities.
Property (P1) says that all node-classification functions fitting sufficiently many generated examples rely on a dependency of length $d$, with arbitrarily high precision and probability.
Property (P2) says that there exists a fitting function expressible by the GNN-based graph learning approaches considered by the paper, including GCNs.
Property (P3) says that the benchmark is fair, because positive and negative examples are produced by the same probabilistic generator \cite{zhou2025glora}.
(P1) addresses path dependence, (P2) avoids an expressivity trap, and (P3) prevents the label from being encoded through an avoidable distribution shift.

The expressibility claim relies on known results about the logical expressiveness of graph neural networks.
The paper invokes results showing that GCNs can express the intended function, and it notes that the other GNN-based systems considered are at least as expressive for the relevant purpose \cite{barcelo2020logical,zhou2025glora}.
The result means that the function is available in principle, not that training will find it.
Thus, when a system fails on GLoRa at length $d$, the absence of a suitable GNN function is not by itself an explanation.
Expressibility does not identify the cause of failure, which may still involve architecture, optimisation, inductive bias, finite data, or the specific path condition.

The proof sketch in Appendix~C gives the intuition behind (P1).
For a fixed $d$, the generator can produce only finitely many graph structures up to the sampled choices in the algorithm.
Each relevant structure has nonzero probability.
The continuous feature values are sampled from a bounded set, and the approximate dependency definition allows a precision tolerance $\delta$.
For any desired probability level, a sufficiently large balanced sample will contain enough matching positive-negative configurations to witness the dependency condition.
In those matched pairs, the examples agree in the ways needed to rule out shortcut functions.
They differ at the designated path features that determine the label.
Any function that fits all these examples must therefore be sensitive to the intended path-aware dependency.

This guarantee is stronger than stating that a long-range rule generated the labels because it quantifies over all fitting functions.
With enough examples, fitting the benchmark forces the function to rely on the dependency rather than merely showing that the intended function exists.
This distinction is the technical core of GLoRa's contribution.

\section{Experimental Protocol}
\label{sec:glora-protocol}

The experiments use GLoRa to test many graph learning systems while varying the generator parameter $d$.
At each value, the protocol keeps the label rule and class balance fixed, while GNN depth is set to let information traverse the sampled source-to-target path.
The accuracy curve therefore estimates where each system stops learning a fitting function under these conditions.

For each generator parameter $d \in \{3,\ldots,15\}$, the paper generates 5000 positive and 5000 negative examples.
Each half is split into training, validation, and test sets with ratio 80:10:10.
The labels are balanced, so random guessing gives accuracy $0.5$.
The benchmark is also noise-free, so a system that has learned a fitting function should approach accuracy $1.0$.
The paper treats accuracy below $0.8$ as evidence that the model has not effectively learned a fitting function for the examples at that length \cite{zhou2025glora}.

The evaluated systems are grouped into four categories that cover ordinary message passing, proposed fixes for known pathologies, and broader interaction mechanisms.
The first category contains vanilla GNN-based systems: GCN, GAT, GraphSAGE, GatedGCN, GGNN, SGC, and GIN.
The second category contains systems designed mainly to address over-smoothing, including GCNII, PairNorm, DropEdge, GPRGNN, DAGNN, APPNP, ResGCN, DenseGCN, JKNet, IGNN, EIGNN, GIND, and related gated or attention variants.
The third category contains systems aimed at over-squashing and related long-range issues, including variants of DRew-GCN, SP-GCN, FA-GCN, DIGL, MixHop-GCN, SDRF, and FOSR.
The fourth category contains transformer-based graph systems: GraphTransformer, SAN, and GraphGPS with positional or structural encodings.
The detailed implementation sources are listed in Appendix~A of the paper \cite{zhou2025glora}.

The training protocol applies a common evaluation framework rather than tuning the benchmark to one architecture.
Each system is run five times on each benchmark.
Training lasts for at most 300 epochs or until convergence.
Hyperparameters are validated, and the reported score is the average test accuracy over the five runs.
For the GNN-based systems, the number of layers is set to $d+2$.
This is enough for the intended source-to-target information to reach the target through message passing, so a failure cannot be attributed to using too few message-passing layers.
The experiments use AdamW, batch normalisation, gradient clipping, ReLU activations where appropriate, and a fixed batch size of 32.
The paper reports that the experiments were run on a shared cluster with one GPU, four CPUs, and up to 32GB of system memory allocated per system \cite{zhou2025glora}.

Appendix~A also reports benchmark statistics for representative lengths.
As the generator parameter $d$ increases, the generated graphs become larger and the average source-to-target path length grows.
For example, the average path length from source to target is reported as 5 for generator parameter $d=5$, 9 for $d=10$, and 14 for $d=15$.
These averages are close to, but not always equal to, the parameter value because the generator varies the exact path length.
The average in-degree remains around 2 in the reported statistics.
This matches the benchmark's design: the task becomes longer-range without turning into an uncontrolled high-degree compression problem.

\section{Main Findings}
\label{sec:glora-main-findings}

The main experiment asks whether the evaluated graph learning systems continue to solve the source-to-target path condition as GLoRa increases the dependency-length parameter.
Across the tested families, the answer is negative at the larger lengths, but the result is specific to the benchmark and protocol.

Many systems fail as $d$ grows, and most show a clear performance drop before $d$ reaches 9.
Some systems perform better than others, but even the best systems perform poorly after $d=11$.
The paper therefore concludes that the evaluated systems do not provide reliable evidence of learning the larger dependency lengths tested by GLoRa \cite{zhou2025glora}.

The best-performing systems use gated edges, especially GatedGCN and gated variants such as DenseGatedGCN and ResGatedGCN.
Gating may help preserve or regulate information as it is propagated along a path.
The result points to a plausible architectural direction, but it does not solve the problem.
Even with gated edges, performance still degrades at modest dependency lengths.

The transformer-based systems do not dominate the benchmark.
Graph transformers change the interaction pattern more radically than ordinary message passing and can, in principle, allow broader interactions between nodes.
Broad attention is not the same as learning the specific source-to-target path condition.
If attention does not learn the right structural relation, the extra interactions do not certify path-aware dependency learning.
GLoRa therefore separates two ideas that are sometimes conflated: the ability to exchange information globally and the ability to learn the intended long-range graph dependency.

The results also clarify what can be concluded about methods designed for over-smoothing and over-squashing.
Such methods can improve deep GNN training and may improve performance on many benchmarks, but on GLoRa they do not automatically produce reliable long-range dependency learning.
This does not mean that the methods are unhelpful.
The benchmark instead tests whether a model that avoids a familiar pathology can learn the path-aware function required by the task.

The accuracy curve supports a length-specific interpretation rather than a single aggregate comparison.
On a GLoRa benchmark with generator parameter $d$, high accuracy would be evidence that the system has learned a fitting function with the dependency guaranteed by Property (P1).
Low accuracy shows that the system did not learn such a function under the experimental protocol.
The curve over increasing $d$ therefore estimates where the system's dependency-length performance begins to fail.

\section{Diagnostic Findings}
\label{sec:glora-diagnostics}

The diagnostic experiments address the ambiguity of low accuracy by testing whether the observed GLoRa failures match the traces predicted by over-smoothing, over-squashing, and vanishing gradients.
They focus on the best-performing systems in three categories: GatedGCN, DenseGatedGCN, and 1DRew-GCN.
The paper does not prove that the three phenomena are irrelevant to long-range graph learning in general.
Rather, the reported traces do not support treating them as explanations for the particular performance drop observed in the inspected GLoRa experiments \cite{zhou2025glora}.

\subsection{Over-smoothing}

If over-smoothing explained the failure, one would expect node representations to converge as the number of layers grows.
The paper therefore compares the largest generator parameter where accuracy is perfect for the selected systems, $d=6$, with the smallest value where accuracy is close to random guessing, $d=13$.
For each case, it extracts the one-dimensional last-layer representation of the target node and plots the distribution across examples.
At $d=6$, the values for positive and negative examples are clearly separated.
At $d=13$, the values are mixed, as expected from poor accuracy.
However, they retain substantial standard deviation rather than collapsing to a single value.
The representations are confused but not smoothed into indistinguishability.
The paper therefore argues that over-smoothing is not the observed reason for failure in this diagnostic.

\subsection{Over-squashing}

Over-squashing concerns the compression of information from many relevant nodes, directed paths, or message-passing walks into one fixed-size representation.
In some graph families, the number of relevant paths or walks into a target can grow exponentially with depth.
The over-squashing diagnostic uses GLoRa's directed construction to check whether this path proliferation occurs.
Each example has a bounded number of additional chains, at most ten in the generator.
Therefore, the number of source-to-target paths remains bounded independently of $d$.
As $d$ increases, the path becomes longer without requiring the target representation to compress an exponentially growing number of paths.

This interpretation depends on directedness.
In an undirected graph, message passing can use both edge directions, and even a simple path can induce many possible message-passing walks into the target.
That would make it harder to separate longer dependency length from the compression of many walks into the target representation.
The directed construction narrows this diagnostic ambiguity because the benchmark structure does not create the many-path bottleneck that over-squashing is meant to describe.

\subsection{Vanishing gradients}

If the failure came from gradients disappearing in early layers, then the first-layer gradients should become very small as the model depth grows.
The vanishing-gradient diagnostic therefore plots the first-layer weight gradients across training epochs for the selected systems.
These gradients remain stable and stay well above zero.
The appendix further reports similar behaviour for other layers.
The measured training dynamics do not support the explanation that gradient signals cannot reach the early layers.

The evaluated systems use standard techniques that help gradient flow, including AdamW, weight initialisation, gradient clipping, ReLU activations, and batch normalisation.
These techniques do not guarantee success on GLoRa, but they weaken the explanation that the models simply could not train because of vanishing gradients.
For the inspected systems, the number of message-passing layers is sufficient and the intended function is expressible.
The measured target representations do not collapse, the directed graphs avoid uncontrolled path compression, and the measured gradients do not vanish.
Yet they still fail at modest dependency lengths.

\section{Thesis-Level Interpretation}
\label{sec:glora-interpretation}

GLoRa's thesis-level role is methodological rather than architectural.
It provides a controlled path-aware test in which shortcut functions, fairness, and expressibility are addressed together.

For each chosen generator parameter $d$, GLoRa creates examples that target a source-to-target path condition and are designed to rule out shortcut functions.
The resulting accuracy curve is evidence about the length of path-aware dependency that a system can learn under these controlled conditions, not only a record of performance on another synthetic dataset.

Empirically, a broad set of state-of-the-art systems fails beyond modest lengths.
Methodologically, this failure occurs on a dataset whose design ties successful fitting to a dependency of specified length while keeping the intended function expressible.
These controls give the negative result a more precise interpretation.

If the diagnostics do not support over-smoothing, over-squashing, or vanishing gradients as explanations for the observed degradation, then the inspected GLoRa failures cannot be reduced to those three labels.
They remain important phenomena, and they may explain failures in other settings.
GLoRa leaves a residual problem: even when the reported traces do not support the usual explanations, current systems may still not learn the intended long-range dependency.
This motivates future work on more precise diagnostics, better inductive biases for path conditions, and architectures that preserve the right kind of information over distance.

The broader implication is that evaluation and modelling are inseparable.
If the benchmark allows shortcut functions, a high score can overstate progress.
If the benchmark is inexpressible, a low score can understate what the model family might learn on a suitable task.
GLoRa addresses both threats by making the dependency explicit, ruling out shortcut functions, and keeping the target within reach of standard GNN-based systems.
This controlled link between the score and the dependency is the contribution developed in the included paper and the reason it forms the core of this thesis.
